Filter alle meldingen die een status hebben van Warning of erger uit de logs van je eigen systeem. Zoek van deze meldingen na wat de reden is dat je deze melding krijgt en wat je er aan zou kunnen doen om deze melding niet meer te zien (geruststelling: dat kan niet altijd).

Maak een tekstdocument met daarin een tabel. In de linker kolom komt de melding, in de middelste kolom de oorzaak en in de rechter kolom de eventuele oplossing.

\begin{itemize}
\item Start Event Viewer (Logboeken)
\item Bekijk de drie belangrijkste categorieën logboeken:
    \begin{itemize}
    \item Windows-logboeken (Application, Security, System, enz.)
    \item Applications and Services Logs
    \item Custom Views
    \end{itemize}
\end{itemize}

Schrijf kort (max. 2 zinnen per onderdeel) waarvoor elk van de Windows-logboeken gebruikt wordt:
\begin{description}
\item[Application] : 
\item[Security] : 
\item[System] : 


3. Analyseren (20 min)
- Zoek in Security-logboek naar inlogpogingen.
- Filter alleen op Event ID 4624 (Successful logon).

Beantwoord de vragen:

Aantal succesvolle logins in de afgelopen 24 uur: ________________________________

Accounts waarop is ingelogd: ________________________________


Extra: Bekijk ook Event ID 4625 (Failed logon). Wat valt je op?

Antwoord: ________________________________

4. Reflectie (15 min)
Beantwoord de vragen:
- Waarom is het filteren van logs belangrijk voor een systeembeheerder?
- Geef twee voorbeelden van situaties waarin je logboeken gebruikt:
  • Voor storingsanalyse
  • Voor beveiliging

Schrijf een kort advies (5–6 zinnen) voor een collega: hoe kun je logboeken efficiënt gebruiken zonder te verdrinken in te veel data?

Antwoord: _________________________________________________________________
